
\documentclass{article}
\usepackage{colortbl}
\usepackage{makecell}
\usepackage{multirow}
\usepackage{supertabular}

\begin{document}

\newcounter{utterance}

\centering \large Interaction Transcript for game `bbh', experiment `hyperbaton', episode 155 with score{-}Qwen3.5{-}9B{-}sp.
\vspace{24pt}

{ \footnotesize  \setcounter{utterance}{1}
\setlength{\tabcolsep}{0pt}
\begin{supertabular}{c@{$\;$}|p{.15\linewidth}@{}p{.15\linewidth}p{.15\linewidth}p{.15\linewidth}p{.15\linewidth}p{.15\linewidth}}
    \# & \multicolumn{2}{c}{Player} && \multicolumn{2}{c}{Game Master} \\
    \hline

    \theutterance \stepcounter{utterance}  
    & & & \multicolumn{4}{p{0.6\linewidth}}{
        \cellcolor[rgb]{0.9,0.9,0.9}{
            \makecell[{{p{\linewidth}}}]{
                \texttt{\tiny{[P1$\langle$GM]}}
                \texttt{Order adjectives correctly in English sentences.} \\
\\ 
\texttt{Q: Which sentence has the correct adjective order:} \\
\texttt{Options:} \\
\texttt{(A) rubber terrible ship} \\
\texttt{(B) terrible rubber ship} \\
\texttt{A: (B)} \\
\\ 
\texttt{Q: Which sentence has the correct adjective order:} \\
\texttt{Options:} \\
\texttt{(A) repulsive small Brazilian exercise ship} \\
\texttt{(B) Brazilian repulsive exercise small ship} \\
\texttt{A: (A)} \\
\\ 
\texttt{Q: Which sentence has the correct adjective order:} \\
\texttt{Options:} \\
\texttt{(A) blue gold wonderful square shoe} \\
\texttt{(B) wonderful square blue gold shoe} \\
\texttt{A: (B)} \\
\\ 
\texttt{Q: Which sentence has the correct adjective order:} \\
\texttt{Options:} \\
\texttt{(A) square cardboard old{-}fashioned Japanese exercise white motorcycle} \\
\texttt{(B) old{-}fashioned square white Japanese cardboard exercise motorcycle} \\
\texttt{A:} \\
            }
        }
    }
     \\ \\

    \theutterance \stepcounter{utterance}  
    & \multicolumn{4}{p{0.6\linewidth}}{
        \cellcolor[rgb]{0.561,0.737,0.561}{
            \makecell[{{p{\linewidth}}}]{
                \texttt{\tiny{[P1$\rangle$GM]}}
                \texttt{(B)} \\
\\ 
\texttt{**Reasoning:**} \\
\texttt{The standard order of adjectives in English follows this sequence:} \\
\texttt{1.  **Opinion** (e.g., *old{-}fashioned*)} \\
\texttt{2.  **Size** (e.g., *square*)} \\
\texttt{3.  **Age** (e.g., *white* {-} though often color, here it functions as a descriptor before material/type in this specific context, or potentially age if interpreted differently, but let's look at the other categories) {-}> Actually, let's re{-}evaluate the specific adjectives:} \\
\texttt{    *   *old{-}fashioned*: Opinion} \\
\texttt{    *   *square*: Size/Shape} \\
\texttt{    *   *white*: Color} \\
\texttt{    *   *Japanese*: Origin/Nationality} \\
\texttt{    *   *cardboard*: Material} \\
\texttt{    *   *exercise*: Purpose/Type (often treated as a noun adjunct or specific type)} \\
\\ 
\texttt{Let's re{-}sort based on the standard mnemonic **OSASCOMP** (Opinion, Size, Age, Shape, Color, Origin, Material, Purpose):} \\
\texttt{1.  **Opinion**: *old{-}fashioned*} \\
\texttt{2.  **Size/Shape**: *square*} \\
\texttt{3.  **Color**: *white*} \\
\texttt{4.  **Origin**: *Japanese*} \\
\texttt{5.  **Material**: *cardboard*} \\
\texttt{6.  **Purpose/Type**: *exercise*} \\
\\ 
\texttt{Order: **old{-}fashioned** (Opinion) {-}> **square** (Shape) {-}> **} \\
            }
        }
    }
    & & \\ \\

    \theutterance \stepcounter{utterance}  
    & & & \multicolumn{2}{p{0.3\linewidth}}{
        \cellcolor[rgb]{0.9,0.9,0.9}{
            \makecell[{{p{\linewidth}}}]{
                \texttt{\tiny{[GM$|$GM]}}
                \texttt{Target: (B)} \\
            }
        }
    }
    & & \\ \\

    \theutterance \stepcounter{utterance}  
    & & & \multicolumn{2}{p{0.3\linewidth}}{
        \cellcolor[rgb]{0.9,0.9,0.9}{
            \makecell[{{p{\linewidth}}}]{
                \texttt{\tiny{[GM$|$GM]}}
                \texttt{game\_result = LOSE} \\
            }
        }
    }
    & & \\ \\

\end{supertabular}
}

\end{document}
